\chapter{Momentum and Trend Following}
\label{ch:momentum}
\chaptermark{Momentum and trend following}

Every chapter so far has bet \emph{against} the last move. The
Roll estimator (\cref{ch:spreads_roll}) infers a bid--ask from
negative autocovariance of trade-to-trade changes; the
Glosten--Milgrom posterior (\cref{ch:kyle_gm}) is a martingale
converging toward fair value as informed and uninformed flow
average out; the OU machinery (\cref{ch:mean_reversion_ou}) trades
the standardised deviation of a stationary process back to its
long-run mean; the pairs and basket spreads
(\cref{ch:pairs_cointegration,ch:basket_etf_arb}) are cointegrating
residuals, traded with the same $z$-score rule. Five chapters, one
instinct: when a price has moved, lean against the move.

This chapter takes the opposite bet: the last move will
\emph{continue}. The signal is the raw past return itself, used as
a forecast of the next return with a positive coefficient. Mean
reversion and momentum are statements about different time scales.
On microsecond horizons, trends are absorbed by market-maker
quoting within milliseconds; on monthly horizons, the opposite is
empirically true. Time series momentum is one of the most robust
anomalies in modern asset pricing, documented across $58$ liquid
futures contracts in four asset classes over two decades by
Moskowitz, Ooi, and Pedersen \citeMOP, and its canonical
formulation fits on a single line. The chapter's job is to make
the coexistence of momentum and mean reversion precise: when
each, and what changes between them.

\begin{backgroundbox}[Prerequisites]
From earlier in the book:
\begin{itemize}
  \item \Cref{ch:mean_reversion_ou}: OU and the mean-reversion
  worldview. Same machinery, sign of the coefficient flipped. A
  return process can have one statistical signature on one horizon
  and the opposite on another.
  \item \Cref{ch:price_value}: the distinction between the
  efficient price (a martingale) and a deviation from slow fair
  value. Momentum is a claim about the efficient price itself, not
  about a spread.
  \item \Cref{ch:participants}: who trades and why. TSMOM's
  behavioural explanation rests on slow-moving allocators (pension
  funds, central banks) whose flow is predictable on monthly
  horizons but invisible at the trade level.
  \item \Cref{ch:order_flow}: Hawkes arrival clustering as the
  microsecond version of trend persistence. ``Momentum'' applies
  to the millisecond and the month, but the mechanism, the signal,
  and the trade differ.
  \item Elementary statistics: standard deviation, annualisation
  ($\sigma_{\text{ann}} = \sigma_{\text{monthly}}\sqrt{12}$), the
  $t$-statistic, sign function.
\end{itemize}
\end{backgroundbox}

\section{Time series momentum (TSMOM)}
\label{sec:ch12_tsmom}

Time series momentum's signal is the simplest non-trivial function
of past prices: the \emph{sign} of the past twelve-month excess
return. Up over the past year: go long. Down: go short. Flat: do
nothing. The position is sized inversely to recent realised
volatility, so a volatile market gets a small position and a calm
one a large position; every contract contributes the same ex-ante
risk. That is the entire strategy.

\begin{definition}[Time series momentum signal]
\index{time series momentum|textbf}\index{momentum!time series|textbf}
\label{def:tsmom_signal}
Let $r_{t-12m,\,t}$ denote the excess return of a tradeable
asset over the twelve months ending at time $t$. The
\emph{time series momentum signal} at time $t$ is
\begin{equation}
  \text{signal}_t \;\defeq\; \operatorname{sign}\!\bigl(r_{t-12m,\,t}\bigr)
  \;\in\; \{-1,\,0,\,+1\},
  \label{eq:tsmom_signal}
\end{equation}
where $\operatorname{sign}(x) = +1$ for $x>0$, $-1$ for $x<0$,
and $0$ for $x = 0$.
\end{definition}

The signal is binary in direction; the position is not. To make
positions comparable across assets, scale notional inversely to a
recent realised volatility estimate $\hat\volat_t$, targeting an
annualised volatility $\volat_{\text{target}}$. \citeMOP{} use
$\volat_{\text{target}} = 40\%$ annualised per instrument, the
level at which a forty-instrument diversified portfolio with
realistic cross-instrument correlations reaches an aggregate target
of about $12\%$ at the strategy level. It is an aggregate-vol
decision projected back onto per-instrument level given
diversification. Their volatility estimator is an exponentially
weighted average of squared daily log-returns with centre-of-mass
about sixty days; any recent realised-vol estimate will do. The
worked example below uses an equally-weighted estimator; treat
the choice as a tuning knob, not a load-bearing assumption.

\begin{keyfactbox}[TSMOM position-sizing formula]
\label{kf:tsmom_formula}
The time series momentum position in a single instrument at
time $t$, in units of notional exposure (positive for long,
negative for short), is
\begin{equation}
  w_t \;=\; \operatorname{sign}\!\bigl(r_{t-12m,\,t}\bigr)
  \cdot \frac{\volat_{\text{target}}}{\hat\volat_t},
  \qquad \volat_{\text{target}} = 0.40 \text{ (annualised).}
  \label{eq:tsmom_position}
\end{equation}
The signal contributes the direction; the volatility ratio
contributes the magnitude. A volatile asset
($\hat\volat_t$ large) gets a small position; a calm asset
gets a large position. The position is rebalanced monthly as
both the signal and the volatility estimate update.
\end{keyfactbox}

First, what the formula does \emph{not} say: it is not ``buy the
asset that went up the most,'' and not a cross-asset ranking. The
signal for gold depends only on gold's own past return; an asset
can be TSMOM-long even if it underperformed every other asset in
the universe. Each instrument's signal is self-contained; the
contrast with cross-sectional momentum is the next section's
subject.

Second, what volatility scaling buys. Without it, one unit of S\&P
futures and one unit of two-year Treasury futures contribute
wildly different risk: the equity index is $5$ to $10$ times
more volatile than the short bond. A naive ``one unit per
long-signal asset'' aggregate would be dominated by the most
volatile contracts, and the portfolio would behave as its equity
sleeve. Inverse-vol scaling forces every asset to contribute equal
ex-ante risk, so diversification works \emph{across} asset
classes.

\subsection{A worked numerical example: gold futures}
\label{sec:ch12_gold_example}

Take gold futures at month-end. Past twelve-month return is
$+15\%$; recent realised volatility is
$\hat\volat_t = 12\%$ annualised. Plug into
\eqref{eq:tsmom_position}:
\begin{equation*}
  w_t
  \;=\;
  \underbrace{\operatorname{sign}(+0.15)}_{=\,+1}
  \;\cdot\;
  \underbrace{\frac{0.40}{0.12}}_{=\,3.333\ldots}
  \;=\;
  +3.33.
\end{equation*}
The position is $+3.33$ units of notional, that is, $3.33$ times
the dollar size of one ``baseline'' unit. The multiplier exceeds one
because gold's $12\%$ volatility is below the $40\%$ target, so
the strategy levers up by $0.40/0.12 \approx 3.333$ to reach
target ex-ante volatility per instrument. At $40\%$ realised vol
the multiplier would be $1$; at $80\%$, $0.5$. A $-15\%$
twelve-month return would flip the sign, giving $-3.33$ (short).

\begin{intuitionbox}
Read TSMOM as ``which way did the slow money push, and how
confident are we given how noisy the asset is?'' The sign of the
past return is the direction slow allocators have collectively
pushed over the past year. Vol-scaling is confidence weighting:
when recent moves are small relative to the trend, the trend
stands out above noise and the bet is bigger; otherwise smaller.
Behavioural story: when a large institution decides to overweight
a sector (say, a sovereign wealth fund adding two percentage
points of NAV to commodities), it does not buy in an hour. It
buys over weeks or months, for impact (\cref{ch:market_impact})
and governance reasons (approval, cash clearing, benchmark
updates). That flow pushes prices in one direction over a long
horizon; TSMOM rides the wave.
\end{intuitionbox}

\subsection{The empirical result}
\label{sec:ch12_empirical}

TSMOM earns its place here on robustness. \citeMOP{} study $58$
liquid futures contracts across four asset classes ($24$
commodity, $9$ equity-index, $13$ bond, $12$ currency forwards)
from January~$1985$ to December~$2009$. Their pooled regression of
next-month vol-scaled return on the sign of the past twelve-month
vol-scaled return,
\begin{equation*}
  \frac{r^s_t}{\sigma^s_{t-1}}
  \;=\; \alpha + \beta_h \,\operatorname{sign}\!\bigl(r^s_{t-h}\bigr)
  + \varepsilon^s_t,
\end{equation*}
yields a positive, statistically significant $\beta_h$ for
$h = 1, 2, \ldots, 12$ months pooled, and within each asset class
individually. The $t$-statistic on $\beta_{12}$ in the all-asset
pooled regression exceeds $5$ with month-clustered standard
errors.

The diversified TSMOM strategy, equal-vol-weighting all $58$
contracts, delivers a gross Sharpe around $1.5$ and monthly alpha
of $1.58\%$ ($t \approx 8.0$) against the Fama--French three-factor
model augmented with the cross-sectional momentum factor
$\textit{UMD}$ of \citet{jegadeesh1993}. The alpha shrinks to
$1.09\%$ ($t \approx 5.4$) with added controls for bond,
commodity, and currency cross-sectional momentum, but survives.
TSMOM is not a relabelling of cross-sectional momentum; it
captures a distinct anomaly. The next section makes the
distinction precise.

\section{Cross-sectional vs time-series momentum}
\label{sec:ch12_cross_vs_ts}

The momentum literature has two parents; do not confuse them.
\citet{jegadeesh1993} introduced \emph{cross-sectional} momentum
(XSMOM) on US equities: rank stocks by past-window return (three
to twelve months), go long the top decile and short the bottom
decile, hold for a similar window. The signal is a \emph{rank}.
The strategy is dollar-neutral and asset-class-neutral by
construction. It works in equities, with variants in country
indices, currencies, and commodities.

Concretely, the two ask different questions. TSMOM asks ``has
\emph{this} asset gone up over the past year?'' If yes, buy
it; if no, short it. XSMOM asks ``has this asset gone up
\emph{more than its peers}?'' Buy the winners, short the
laggards, regardless of absolute direction. In a year when every
stock is up $20\%$, TSMOM is long everything; XSMOM is long the
top-decile $+40\%$ names and short the bottom-decile $+5\%$ names
for zero net exposure.

The younger parent is \citeMOP ($2012$) time series momentum. The
signal is each asset's own past return, not its rank: the strategy
can simultaneously be long every asset (if all $58$ are up on the
year) or short every asset (if all are down). There is no
mechanical dollar- or asset-class-neutrality. Inverse-vol scaling
equalises risk across instruments, but net market exposure is
whatever the signs collectively dictate.

\begin{keyfactbox}[Two kinds of momentum]
\label{kf:two_momentums}
\renewcommand{\arraystretch}{1.2}
\begin{center}
\begin{tabular}{p{0.27\linewidth} p{0.31\linewidth} p{0.31\linewidth}}
\toprule
 & \textbf{Cross-sectional (XSMOM)} & \textbf{Time series (TSMOM)} \\
\midrule
Signal &
$\operatorname{rank}\!\bigl(r^s_{t-12,\,t}\bigr)$
within universe &
$\operatorname{sign}\!\bigl(r^s_{t-12,\,t}\bigr)$
of own return \\
Positions &
Long top decile, short bottom decile (dollar neutral) &
Long if own past return $> 0$, short if $< 0$ (any net) \\
Net exposure &
Zero by construction &
Time-varying; can be net long, net short, or zero \\
Canonical paper &
Jegadeesh--Titman (1993) on US equities &
Moskowitz--Ooi--Pedersen (2012) on $58$ futures \\
Universe &
Often single-asset-class (e.g.\ equities) &
Multi-asset-class futures \\
Captured by &
Fama--French $\textit{UMD}$ factor &
\emph{Not} fully captured by $\textit{UMD}$; has additional alpha \\
\bottomrule
\end{tabular}
\end{center}
\end{keyfactbox}

The two are positively correlated (if equities rally across
the board, both XSMOM and TSMOM tend to be long equities) but
not identical. \citeMOP{} regress TSMOM returns on $\textit{UMD}$
and find significant positive loading with residual alpha around
$1.09\%$ per month ($t > 5$). Even after subtracting what
cross-sectional momentum explains, TSMOM still earns money;
time-series direction adds incremental information.

Three operational consequences. First, Prosperity's typical
universe is a handful of products with no sensible cross-section,
so only the time-series formulation translates. Second, a Goldman
macro desk's trend pod is multi-asset (currencies, rates,
equities, commodities); even with XSMOM-style ranking within an
asset class, the aggregate portfolio combines both kinds of
momentum. Third, the two effects are independent, so blending
yields a small Sharpe boost over either alone.

\section{Why momentum works on monthly horizons but not microseconds}
\label{sec:ch12_horizons}

Five chapters have argued mean reversion; this one argues the
opposite. Both cannot hold on the same time scale for the same
object, so what is going on?

They are true on \emph{different} time scales, with a different
mechanism dominating each. Time scale, not asset class, is the
key axis. From fastest to slowest:

\paragraph{Microsecond horizon ($10^{-6}$ to $10^{-3}$ seconds).}
A market-making algorithm sees an aggressive buy, infers a small
adverse-selection signal, widens its ask and lifts its bid
(\cref{ch:market_making}, Avellaneda--Stoikov skew). The next
aggressive buy meets a higher quote; within milliseconds the
quoting algorithm reverts price toward fair value. Market-makers
\emph{eat} momentum on this horizon: any short-burst trend in
trade prices is absorbed into the spread within milliseconds. The
Roll spread of \cref{ch:spreads_roll} is this effect viewed
through autocovariance: $\Cov(\Delta p_t, \Delta p_{t+1}) = -s^2/4
< 0$; trade-level returns are \emph{negatively} autocorrelated.

\paragraph{Sub-second to one-minute horizon ($10^{-3}$ to $60$
seconds).} Mixed regime. Hawkes-style arrival clustering
(\cref{ch:order_flow}) creates positive autocorrelation in
\emph{order flow}, and imbalance predicts the next-second mid move
with a small positive coefficient. It is the closest thing to
``momentum'' in the millisecond--second band, but the prediction
is about \emph{order flow} continuing, not trends. The signal
decays within seconds and is invisible on monthly bars.

\paragraph{Minutes to hours.} Genuinely mixed. Some markets trend
because execution algorithms break parent orders into child orders
(the \cref{ch:almgren_chriss} machinery spreads trades over an
hour, generating predictable drift in the execution window);
others mean-revert under contrarian market-maker inventory
pressure. Tradeability depends on liquidity profile and
participant mix.

\paragraph{Days to weeks.} Mixed but momentum-leaning for liquid
macro futures. The MOP regressions give positive coefficients down
to one-month lookbacks; Chan \citeChan{} (Ch~$6$) reports
practitioner experience with $1$--$5$ day momentum on FX carry
pairs and commodity futures. Equity pairs and ETF baskets
mean-revert on the same horizon because the dominant force is
cointegration convergence. \Cref{ch:which_strategy_when} diagnoses
which regime a product is in; the rule of thumb is that single
instruments with macro drivers trend, and spreads of related
instruments revert.

\paragraph{One to twelve months.} Strongly momentum-favouring
across all four asset classes \citeMOP. Three non-exclusive
explanations coexist in the literature: (i)~\emph{underreaction}
to gradual fundamental news (a slow inflation trend, a
multi-quarter earnings-revision drift) means prices move in
steps rather than jumps; (ii)~\emph{investor herding} and
delegated-manager career concerns draw late capital toward
already-moving assets; (iii)~\emph{slow capital allocation}
(sovereign funds, pension plans, insurance asset-liability desks)
moves on monthly-to-quarterly cycles, and once a flow direction
is established it persists because allocator decisions are
rebalanced infrequently. All three produce the same statistical
signature (positive autocorrelation of monthly returns), and
all three decay: the edge is strong over weeks to months but
flips sign beyond about a year, as the next paragraph shows.

\paragraph{Beyond one year.} The MOP regressions show
\emph{negative} coefficients beyond about thirteen months:
long-horizon mean reversion that DeBondt and Thaler documented for
equities in the $1980$s and that has been reproduced for
commodities and currencies. The five-year reversal reads as slow
correction of multi-year mispricing, as opposed to slow
propagation of new information driving $1$--$12$ month momentum.

\begin{pitfallbox}[Don't run monthly momentum on a five-minute bar]
\label{pf:wrong_horizon}
The \citeMOP{} edge is about \emph{monthly} horizons on
\emph{liquid macro futures}. It does not translate to
higher-frequency data without substantial modification:
\begin{itemize}
  \item At a five-minute lookback the sign-of-past-return signal
  is dominated by microstructure noise. Roll and
  Glosten--Milgrom both imply short-horizon \emph{negative}
  autocorrelation at the trade level. Five-minute TSMOM buys
  every uptick, sells every downtick, pays spread on every flip,
  and loses to market-makers.
  \item The $40\%$ annualised target was calibrated for monthly
  rebalancing on contracts with single-digit monthly volatilities.
  Five-minute realised vols give nonsense leverage: the
  five-minute-to-annualised conversion factor is an order of
  magnitude larger than the monthly one
  ($\sqrt{78 \cdot 252} \approx 140$ on a $6.5$-hour equity
  session, $\sqrt{288 \cdot 252} \approx 269$ round-the-clock).
  Comparing scaled five-minute dispersion to
  $\volat_{\text{target}} = 0.40$ gives ratios unrelated to the
  investment-horizon risk the formula controls.
  \item The behavioural mechanism, slow allocator flow, has
  no analogue at the second or minute horizon. High-frequency
  trend persistence comes from order-flow clustering
  (\cref{ch:order_flow}) or parent-order execution drift
  (\cref{ch:almgren_chriss}), not allocator behaviour.
\end{itemize}
Rule: pick the model whose mechanism matches the horizon. Hawkes
for milliseconds, Roll for trade-level, OU for spreads on hours
and days, TSMOM for macro futures on weeks and months, five-year
reversal beyond a year.
\end{pitfallbox}

\section{Momentum crashes}
\label{sec:ch12_crashes}

TSMOM works on average over long samples, but its return
distribution has a heavy left tail. It is occasionally the
worst-positioned strategy in the world, exactly at regime turns;
the $2009$ crash is the textbook case.

The mechanism is mechanical. Momentum is long what has gone up,
short what has gone down. After a sustained bear market, the longs
concentrate in defensives (bonds, gold, defensive sectors) and the
shorts in cyclicals (small caps, banks, materials). When the
market bottoms and reverses violently, as in March $2009$,
the cyclical shorts rally hardest and the defensive longs lag. The
portfolio is short the rally and long the laggards; both legs
lose.

Momentum is \emph{short gamma} with respect to regime turns. A
gamma-positive position (a long option) profits from large moves
in either direction; gamma-negative loses. Momentum is
gamma-negative because it bets the recent direction continues:
any large reversal moves the underlying against the position and
forces a rebalance; both are losses. \citet{daniel2016momentum}
formalise this as ``momentum is short a call option on the market
when the market has been down, and short a put option when it has
been up''.

\begin{historybox}
\label{hist:march2009}
The March $2009$ momentum crash is the canonical failure mode.
The S\&P~$500$ bottomed at $666$ on $6$~March~$2009$ after a $57\%$
peak-to-trough drop from the October $2007$ high; by end-April it
had rallied about $30\%$. Cross-sectional momentum in US equities
(short banks and cyclicals, long defensives and gold) lost
$25$--$30\%$ of NAV in March and April alone. AQR Capital's
flagship momentum fund (Cliff Asness) took a comparable drawdown.
This single episode explains a meaningful fraction of post-crisis
underperformance of academic momentum factors. Cross-asset TSMOM fared better than
equity-only XSMOM that quarter (bond and currency sleeves
diversified) but still recorded one of its worst quarters in
the MOP sample. Live trend-following CTAs reported single-month
drawdowns of $5$--$15\%$ and several prominent funds closed. The
industry lesson is that \emph{trend-following needs a tail-risk
overlay}: stop-outs, cross-asset diversification, and (at the
largest funds) options-based protection against violent reversals.
\end{historybox}

Practical mitigations fall into three buckets. \emph{Cross-asset
diversification} is cheapest: bond and currency legs do not crash
on the same days as equities, so a diversified $58$-instrument
portfolio has shallower drawdowns. \emph{Strategy-level vol
targeting} (added on top of per-instrument targeting in the
position formula) cuts exposure when realised vol spikes, which
approximates cutting exposure near regime turns. \emph{Explicit
tail hedges} using OTM options on major equity indices are
expensive in calm periods but pay off precisely when momentum
crashes; the largest trend funds run a permanent options overlay.
Option mechanics in \cref{ch:options_payoffs,ch:greeks_hedging};
Kelly sizing of tail-protected portfolios in \cref{ch:kelly_risk}.

\begin{prosperitybox}
\label{prosp:tsmom_squidink}
The cleanest Prosperity TSMOM application is \texttt{SQUID\_INK}
in Round~$1$. The product was officially mean-reverting
short-term, but the Frankfurt Hedgehogs'
inspection (\texttt{imc-prosperity-3/README.md}, Round~$1$) found
an anonymous bot, later identified as Olivia, consistently buying
$15$ lots at the daily low and selling $15$ at the daily high.
Whenever a trade printed at a daily extreme in a direction
consistent with Olivia's pattern, the Hedgehogs flagged it as a
directional signal and held to the opposing extreme. Their
\texttt{InkTrader} class
(\texttt{FrankfurtHedgehogs\_polished.py}, lines $382$--$405$)
implements this: when \texttt{informed\_direction == LONG}, set
\texttt{expected\_position = position\_limit} and bid the ask wall
toward the cap; symmetrically for SHORT.

This is a discrete cousin of TSMOM. The signal is not the sign of
the past twelve-month return (twelve months exceeds a Prosperity
round) but the sign of the inferred informed-trader direction, a
proxy for ``which way is the slow money pushing''. Position sizing is
hard-clipped to the limit rather than vol-scaled, a degenerate
form of the same ``size-to-target-risk'' instinct. Once the trader
is identified directly via trade ID in Round~$5$, false positives
drop to near zero and the strategy becomes the dominant
SQUID\_INK PnL contributor. A pure TSMOM with a five-tick lookback
without the Olivia overlay yields only a modest edge
(noise-to-signal at five ticks is too high and the product only
marginally trending), but combined with the informed-trader
detector the same time-series-direction instinct produced a
clean, high-Sharpe contributor that depended on detecting slow
capital flow, not predicting it. The overall pattern matches this
chapter's trend-following theme: small but consistent, with upside
from disciplined risk control rather than clever signal
engineering.
\end{prosperitybox}

\begin{gsbox}
\label{gs:tsmom_macro_desk}
On a Goldman systematic macro desk, TSMOM is one ingredient in a
``trend pod'' running on $50$--$100$ liquid futures and forwards
(G$10$ FX, on-the-run government bonds, major equity indices,
energy and metals). The Strat owns the signal pipeline:
return-windowing on adjusted-rolled futures series, the vol
estimator (EWMA with a $60$-day half-life), the $40\%$ target
scaling, and the monthly rebalance schedule. The desk trader owns
execution, which on a monthly rebalance with hundreds of positions
is a small Almgren--Chriss problem (\cref{ch:almgren_chriss}) per
name: the Strat estimates the target, the trader spreads the
trade across hours to manage impact, and post-trade attribution
reconciles realised against model entry. Risk is monitored
separately: strategy-level VaR, gross and net leverage, pairwise
correlation with peer CTAs (Man AHL, Winton, Aspect, Lynx). When
all trend funds are positioned identically, the risk team imposes
a temporary gross cap against crowded-trade unwinds. Drawdowns are
correlated across managers, so the primary defence is not
within-strategy diversification (which trend cannot supply through
a regime turn) but position limits and options overlays on the
equity-index sleeve, sized via the gamma-and-vega framework of
\cref{ch:greeks_hedging}.
\end{gsbox}

\section{A short reference implementation}
\label{sec:ch12_implementation}

The TSMOM signal fits in fewer than twenty lines of Python. The
function below takes a 1-D array of historical prices (most recent
last), the momentum lookback in months, the vol lookback, and the
annualised target, and returns a signed notional position. No
state, no portfolio code, no backtester; those belong in
\cref{ch:backtesting}. The point is to make
\eqref{eq:tsmom_position} unambiguous and reproducible.

\begin{lstlisting}[caption={Time series momentum signal with vol scaling.}]
import numpy as np
def tsmom_signal(prices: np.ndarray,
                 lookback_months: int = 12,
                 vol_lookback: int = 60,
                 vol_target: float = 0.40) -> float:
    """Time series momentum signal with vol scaling.
    prices: array of monthly prices, most recent last.
    Returns position size in units of notional, signed.
    """
    if len(prices) < max(lookback_months + 1, vol_lookback + 1):
        return 0.0
    past_ret = prices[-1] / prices[-lookback_months - 1] - 1.0
    sign = 1.0 if past_ret > 0 else (-1.0 if past_ret < 0 else 0.0)
    log_rets = np.diff(np.log(prices[-vol_lookback - 1:]))
    vol_ann = log_rets.std(ddof=1) * np.sqrt(12.0)
    if vol_ann < 1e-8:
        return 0.0
    return sign * (vol_target / vol_ann)
\end{lstlisting}

Three production notes not encoded in the snippet. First, the vol
estimator is an equally-weighted standard deviation of the last
\texttt{vol\_lookback} monthly log-returns; MOP use an EWMA with
$60$-day-equivalent centre of mass computed on \emph{daily}
returns and annualised. EWMA reacts faster to vol regime shifts,
but the strategy is qualitatively robust. Second, the function is
stateless and assumes a clean monthly price series; in practice
the trader handles futures roll adjustments, missing data, and
timing conventions. Third, the returned position is notional in
an unspecified unit; the caller translates ``notional $+3.33$''
into contracts given the multiplier and current price. None of
these production concerns changes the core signal of
\eqref{eq:tsmom_position}.

\section{Key facts to memorise}
\label{sec:ch12_keyfacts}

\begin{itemize}
  \item \textbf{TSMOM signal.} The sign of the past twelve-month
  excess return predicts the next-month return with a positive
  coefficient. Documented across $58$ liquid futures in four
  asset classes over $1985$--$2009$ by \citeMOP.
  \item \textbf{Position formula.}
  $w_t = \operatorname{sign}(r_{t-12m,t}) \cdot
  (\volat_{\text{target}}/\hat\volat_t)$ with
  $\volat_{\text{target}} = 40\%$ annualised per instrument.
  Direction from sign; magnitude from inverse-vol.
  \item \textbf{Vol-scaling rationale.} Equalises ex-ante risk
  across instruments so a diversified portfolio is not dominated
  by the most volatile asset.
  \item \textbf{XSMOM vs TSMOM.} XSMOM ranks within a universe
  (dollar- and asset-class-neutral). TSMOM judges each asset
  against its own history (time-varying net exposure). Positively
  correlated but not identical; TSMOM retains alpha after
  controlling for $\textit{UMD}$.
  \item \textbf{Horizon beats asset class.} Momentum monthly, mean
  reversion at milliseconds. Behavioural underreaction and slow
  capital drive monthly momentum; adverse selection and
  Glosten--Milgrom convergence drive trade-level reversion. Match
  model to horizon.
  \item \textbf{Crash risk.} Momentum is short gamma with respect
  to regime turns: a violent reversal loses on both legs. March
  $2009$: equity momentum lost $25$--$30\%$ in two months as
  cyclical shorts rallied and defensive longs lagged.
  \item \textbf{Mitigations.} Cross-asset diversification
  (cheapest), strategy-level vol targeting, and options-based tail
  hedges (most expensive, most effective in a crash).
  \item \textbf{Operational rule.} The MOP edge is monthly bars on
  liquid macro futures. It does \emph{not} translate to
  five-minute bars without reformulation; intraday ``trend''
  machinery is order-flow clustering (\cref{ch:order_flow}), not
  TSMOM.
\end{itemize}
